\section{Elaborate}

\paragraph{Input and output.}
Elaboration consumes a normalized closed AST and produces a closed AST with
implicit call structure made explicit where the language contract requires it.

\paragraph{Contract.}
The current elaboration stage is signature-aware and focuses on the boundary
between surface calls and the flat core.  It preserves user-visible type
ascriptions and does not erase inferred metadata.  Its job is to make
defaultable call structure explicit enough that flattening, typecheck2, and
Graph IR lowering do not need ad hoc fallback logic for omitted optional
parameters.

Elaboration accounts for already-supplied positional arguments, explicit
keyword arguments, path arguments introduced by normalization, and optional
parameters with implicit \code{null} defaults.  It adds a missing keyword
argument only when the corresponding formal parameter is not already covered by
the call.  This is part of the round-trip contract: flat rendered Axon
represents scoped calls with ordinary leading \code{Path} arguments, and
reparsing such a program preserves the same call surface instead of adding
duplicate defaults for parameters supplied positionally.

\paragraph{Rationale.}
Elaboration sits after normalization because normalization first determines
which expressions are calls.  Once calls are explicit, elaboration can apply
signature-aware rules without reinterpreting syntax.

\paragraph{Formal view.}
For normalized AST \(A_n\), elaboration produces \(A_e\):
\[
  \mathsf{elaborate}: A_n \rightarrow A_e .
\]
The transformation is signature-aware but not type-inferential.  For a call
\(\code{f}\ a_1 \ldots a_m\ \{k_j=v_j\}\), it binds actual arguments to the
callee surface after accounting for path slots and positional parameters.  For
each remaining formal parameter \(p\), elaboration may insert \(p=d_p\) when
\(p\) has a declared default \(d_p\), or \(p=\code{null}\) when \(p\) is
optional without an explicit default.  It does not solve tensor shapes.

The stage is idempotent under the weak and strong round-trip tests used for
model Axon files: rendering and reparsing an elaborated artifact, with or
without re-running earlier resolution, yields the same elaborated call
structure.

\paragraph{Example.}
If a library function has the signature
\begin{verbatim}
linear :: Path -> Tensor[B,S,D] -> ?Dim -> ?Bool -> Tensor[B,S,O]
linear path x ?dim=null ?bias=true = ...
\end{verbatim}
then a call such as
\begin{verbatim}
y <- NN.linear @w x
\end{verbatim}
can be elaborated to an equivalent call with explicit defaulted arguments:
\begin{verbatim}
y <- NN.linear @w x dim=null bias=true
\end{verbatim}
This prevents downstream stages from inventing backend-specific defaults.
However, if a call already supplies the optional parameter positionally,
elaboration does not also add the keyword form:
\begin{verbatim}
y <- f path x supplied_optional
\end{verbatim}
must remain a single binding of the optional formal, not become
\begin{verbatim}
y <- f path x supplied_optional optional=null
\end{verbatim}
after rendering and reparsing.
